\chapter{Tiling, Ordering and Compositing at Scale}\label{04-visual-computing-ii-ch039-tiling-ordering-and-compositing-at-scale}\phantomsection\hypertarget{04-visual-computing-ii-ch039}{}\label{stagechapter:04-visual-computing-ii-ch039}

\hypertarget{04-visual-computing-ii-ch039-tiling-ordering-and-compositing-at-scale}{}

\hyperref[volume-reading-map]{Reading map} · \hyperref[04-visual-computing-ii-ch038-projecting-ellipsoids-into-screen-space-conics]{Previous chapter} · \href{../04-Light-Transport-and-Image-Synthesis-P01/main.pdf\#nameddest=04-visual-computing-ii-ch040-lighting-reconstruction-and-surface-motion}{Next chapter (Light Transport and Image Synthesis, Part 1, Ch. 10)}

Many projected kernels must be assigned to tiles, ordered and composited efficiently. This span connects local footprints to scalable rasterization and identity-aware data flow.

\textbf{By the end:} Describe the tiling, sorting and compositing stages for large kernel sets.

\textbf{Recall question:} Why does ordering matter when overlapping kernels are composited?

\section{Rasterize kernels at scale and connect their identities}\label{04-visual-computing-ii-ch039-rasterize-kernels-at-scale-and-connect-their-identities}

\hypertarget{04-visual-computing-ii-ch039-rasterize-kernels-at-scale-and-connect-their-identities}{}

\textbf{Policy distinction:} the source opacity cutoff \(A>0.999\) means remaining transmittance below \(10^{-3}\). The inspected Graphdeco reference uses \(10^{-4}\); these are separate accuracy/performance policies, as derived above.

\subsection{+ M150}\label{04-visual-computing-ii-ch039-m150}

\hypertarget{04-visual-computing-ii-ch039-m150}{}

This pass completes \textbf{Phase 25}. The authoritative curriculum boundary is exact: \textbf{M148 owns Parts 2453--2474}, taking the projected conics from Pass 15 through tiling, prefix sums, reference expansion, ordering and alpha compositing; \textbf{M149 owns 2475--2486}, covering the CPU oracle, Metal implementation and quantitative CPU↔GPU validation; and \textbf{M150 owns 2487--2500}, where Gaussian splats stop being merely pixels and become an editable captured-world representation connected to persistent entities and physical degrees of freedom.

Pass 15 ended with one visible Gaussian represented approximately as:

\[
G_i^{screen}
=
\{
\mu_i^{2D},
z_i,
C_i=\Sigma_{2D,i}^{-1},
r_i,
\alpha_i,
c_i,
ID_i
\}.
\]

That representation answers:

\begin{quote}
Where is the Gaussian and what is its elliptical pixel-space kernel?
\end{quote}

It does \textbf{not} yet answer:

\begin{quote}
Which of millions of Gaussians should every pixel examine, and in what order?
\end{quote}

That is the problem M148 solves.

The complete high-performance pipeline becomes:

\[
\boxed{
\text{3D Gaussians}
\rightarrow
\text{project/cull}
\rightarrow
\text{tile-overlap counts}
\rightarrow
\text{exclusive scan}
\rightarrow
\text{tile references}
\rightarrow
\text{stable tile+depth sort}
\rightarrow
\text{tile ranges}
\rightarrow
\text{per-pixel front-to-back composite}
}
\]

Maveb currently implements essentially this entire chain in Metal.

\begin{center}\rule{0.5\linewidth}{0.5pt}\end{center}

\subsection{TILING, REFERENCE EXPANSION, SORTING AND COMPOSITING}\label{04-visual-computing-ii-ch039-tiling-reference-expansion-sorting-and-compositing}

\hypertarget{04-visual-computing-ii-ch039-tiling-reference-expansion-sorting-and-compositing}{}

Imagine:

\[
N=5,000,000
\]

Gaussians and:

\[
1920\times1080
\approx
2.07\times10^6
\]

pixels.

The absolutely naive algorithm is:

\begin{verbatim}
for every pixel:
    for every Gaussian:
        test Gaussian
\end{verbatim}

Cost:

\[
O(PN).
\]

Numerically:

\[
2.07\times10^6
\times
5\times10^6
\approx
10^{13}
\]

Gaussian-pixel tests.

Impossible for interactive rendering.

But each projected Gaussian typically affects only a small screen region.

So instead partition the screen into fixed tiles.

Maveb uses:

\[
\boxed{16\times16}
\]

pixel tiles.

\begin{center}\rule{0.5\linewidth}{0.5pt}\end{center}

\textbf{Tile-grid dimensions}

For image:

\[
W\times H
\]

and tile side:

\[
T,
\]

columns:

\[
\boxed{
N_x
=
\left\lceil
\frac WT
\right\rceil
=
\frac{W+T-1}{T}
}
\]

with integer division.

Rows:

\[
\boxed{
N_y
=
\frac{H+T-1}{T}.
}
\]

Total:

\[
N_{tiles}=N_xN_y.
\]

Maveb does:

\begin{verbatim}
tileColumns =
    (width + 15) / 16;

tileRows =
    (height + 15) / 16;
\end{verbatim}

because:

\[
T=16.
\]

For 1920×1080:

\[
N_x=120,
\]

\[
N_y=
\lceil1080/16\rceil
=
68.
\]

So:

\[
N_{tiles}=8160.
\]

Much smaller than pixel count.

\begin{center}\rule{0.5\linewidth}{0.5pt}\end{center}

\textbf{Why 16×16?}

There is no universal magical tile size.

Tradeoff:

Small tile:

\[
\text{fewer irrelevant Gaussians per pixel}
\]

but:

\[
\text{more tile references and scheduling metadata}.
\]

Large tile:

\[
\text{fewer references}
\]

but every pixel may iterate more irrelevant splats.

For projected Gaussian rasterization, 16×16 is a strong practical compromise and is also the tile size hardcoded in current Maveb:

\begin{verbatim}
constant uint
aetherGaussianTileSize = 16;
\end{verbatim}

\begin{center}\rule{0.5\linewidth}{0.5pt}\end{center}

From Pass 15 we have screen-space center:

\[
(c_x,c_y)
\]

and conservative radius:

\[
r.
\]

Pixel bounds:

\[
x_{min}
=
\max(
0,
\lfloor c_x-r\rfloor
),
\]

\[
x_{max}
=
\min(
W-1,
\lceil c_x+r\rceil
).
\]

Likewise y.

Then tile coordinates:

\[
\boxed{
t_{x,min}
=
\left\lfloor
\frac{x_{min}}T
\right\rfloor
}
\]

and:

\[
\boxed{
t_{x,max}
=
\left\lfloor
\frac{x_{max}}T
\right\rfloor.
}
\]

Same for y.

Maveb computes exactly:

\begin{verbatim}
minimumTile =
    uint2(minimumPixel) / 16;

maximumTile =
    uint2(maximumPixel) / 16;
\end{verbatim}

\begin{center}\rule{0.5\linewidth}{0.5pt}\end{center}

\textbf{Number of overlapped tiles}

If inclusive tile range:

\[
[t_{x0},t_{x1}]
\times
[t_{y0},t_{y1}],
\]

number of tiles touched:

\[
\boxed{
K_i
=
(t_{x1}-t_{x0}+1)
(t_{y1}-t_{y0}+1).
}
\]

Maveb stores this as:

\begin{verbatim}
overlap
\end{verbatim}

and writes:

\begin{verbatim}
tileCounts[index] =
    overlap;
\end{verbatim}

Important nuance:

Despite the name \texttt{tileCounts}, this buffer is initially indexed by \textbf{Gaussian}.

So:

\[
tileCounts[i]
=
\text{number of tile references Gaussian }i\text{ must generate}.
\]

It is not initially:

\[
\text{number of Gaussians in tile }t.
\]

That distinction matters.

\begin{center}\rule{0.5\linewidth}{0.5pt}\end{center}

Suppose:

\[
G_0
\]

touches 4 tiles,

\[
G_1
\]

touches 2,

\[
G_2
\]

touches 5.

Then reference counts:

\[
K=
[4,2,5].
\]

Total flattened references required:

\[
\boxed{
M=
\sum_i K_i
=
11.
}
\]

Why references rather than storing each Gaussian once?

Because one Gaussian may need to appear in several tiles.

The relation is many-to-many:

\[
Gaussian_i
\leftrightarrow
Tile_j.
\]

A Gaussian touching:

\[
20
\]

tiles must appear in:

\[
20
\]

tile work lists.

\begin{center}\rule{0.5\linewidth}{0.5pt}\end{center}

\textbf{Direct tile counting versus Maveb's Gaussian counting}

There are two broad constructions.

A traditional tile-centric approach first counts:

\[
count[t]
=
\text{number of Gaussians touching tile }t.
\]

Then scan tile counts.

Maveb instead counts:

\[
count[g]
=
\text{number of tiles touched by Gaussian }g.
\]

Then scans Gaussian counts to allocate one contiguous output region per Gaussian.

After all references are generated, it sorts them by tile and depth.

Both approaches ultimately produce:

\[
\text{tile}\rightarrow
\text{ordered Gaussian references}.
\]

Maveb's approach is especially natural because every projection thread already knows the entire tile rectangle of its Gaussian.

\begin{center}\rule{0.5\linewidth}{0.5pt}\end{center}

We now know how many references each Gaussian emits:

\[
K_0,K_1,\ldots,K_{N-1}.
\]

We need non-overlapping write ranges.

Exclusive prefix sum:

\[
O_0=0
\]

and:

\[
\boxed{
O_i=
\sum_{j=0}^{i-1}K_j.
}
\]

Then Gaussian \(i\) writes references beginning at:

\[
O_i.
\]

Example:

\[
K=
[4,2,5,1].
\]

Exclusive scan:

\[
O=
[0,4,6,11].
\]

Writes become:

\[
G_0:[0,4)
\]

\[
G_1:[4,6)
\]

\[
G_2:[6,11)
\]

\[
G_3:[11,12).
\]

No atomic increment is needed for every emitted reference.

That is why prefix sum is such a central GPU primitive.

\begin{center}\rule{0.5\linewidth}{0.5pt}\end{center}

\textbf{Why not one global atomic counter?}

Could do:

\begin{verbatim}
uint destination =
    atomic_fetch_add(
        globalCounter,
        overlap);
\end{verbatim}

once per Gaussian.

That can work.

But a scan gives deterministic offsets and avoids potentially severe contention around a global allocator.

More generally:

\[
\boxed{
\text{count}
\rightarrow
\text{scan}
\rightarrow
\text{scatter}
}
\]

is one of the canonical parallel-GPU patterns.

You already learned this in Phase 22.

Here it becomes a real renderer.

\begin{center}\rule{0.5\linewidth}{0.5pt}\end{center}

\textbf{Maveb's parallel scan}

Current default Metal path uses 256-thread blocks.

Each workgroup loads:

\[
256
\]

counts into:

\begin{verbatim}
threadgroup uint scratch[256];
\end{verbatim}

Then performs a Blelloch-style exclusive scan.

Up-sweep:

\[
1,2,4,8,\ldots
\]

combines partial sums.

At the root:

\[
scratch[255]
\]

contains block total.

That total is saved to:

\begin{verbatim}
blockSums[groupIndex].
\end{verbatim}

Root then becomes:

\[
0.
\]

Down-sweep propagates exclusive prefixes.

Finally:

\[
offsets[index]
=
scratch[localIndex].
\]

\begin{center}\rule{0.5\linewidth}{0.5pt}\end{center}

\textbf{Why Blelloch has two phases}

Suppose values:

\[
[3,1,7,0,4,1,6,3].
\]

Up-sweep reduces tree toward total:

\[
25.
\]

But that alone does not provide every prefix.

Set root to zero.

Down-sweep distributes sum of all elements preceding each node.

Result:

\[
\boxed{
[0,3,4,11,11,15,16,22].
}
\]

That is an exclusive scan.

\begin{center}\rule{0.5\linewidth}{0.5pt}\end{center}

\textbf{Multi-block scan}

If Gaussian count exceeds 256, each workgroup gives only local prefixes.

Suppose block \(b\) total:

\[
S_b.
\]

You need prefix of block totals:

\[
B_b
=
\sum_{k<b}S_k.
\]

Then each local offset:

\[
\boxed{
O_i=
O_i^{local}
+
B_{block(i)}.
}
\]

Maveb pipeline therefore executes:

\protect\phantomsection\label{4z8tmw}%
\begin{verbatim}
Gaussian Block Scan
        ↓
Gaussian Block Prefix
        ↓
Gaussian Block Offset Propagation
\end{verbatim}

This is exactly the hierarchical scan algorithm.

\begin{center}\rule{0.5\linewidth}{0.5pt}\end{center}

\textbf{Maveb has a deterministic serial fallback}

If hardware cannot provide required:

\[
256
\]

threads per threadgroup, Maveb keeps:

\protect\phantomsection\label{81ri4m}%
\begin{verbatim}
Gaussian Serial Scan Fallback.
\end{verbatim}

This is an important engineering principle:

\[
\boxed{
\text{correctness path remains executable independently of optimized path}.
}
\]

Its ADR explicitly preserves a bounded fallback while performance claims must use the parallel path.

\begin{center}\rule{0.5\linewidth}{0.5pt}\end{center}

CSR normally means Compressed Sparse Row.

For sparse matrix:

\protect\phantomsection\label{jurd6v}%
\begin{verbatim}
rowOffsets
columnIndices
\end{verbatim}

For Gaussian tiles, analogous structure is:

\protect\phantomsection\label{re9vks}%
\begin{verbatim}
tileRanges
sortedGaussianIndices
\end{verbatim}

For tile \(t\):

\[
begin=
ranges[t].x
\]

\[
end=
ranges[t].y.
\]

Its splats are:

\[
sortedValues[
begin:end
].
\]

That is essentially CSR.

\begin{center}\rule{0.5\linewidth}{0.5pt}\end{center}

\textbf{Example}

Suppose final sorted tile references are:

\protect\phantomsection\label{usz654}%
\begin{verbatim}
Tile 0 → G2 G7 G9
Tile 1 → G1
Tile 2 → G0 G4
\end{verbatim}

Flatten:

\protect\phantomsection\label{m5jiq5}%
\begin{verbatim}
values =
[
 G2,G7,G9,
 G1,
 G0,G4
]
\end{verbatim}

Ranges:

\[
range_0=[0,3),
\]

\[
range_1=[3,4),
\]

\[
range_2=[4,6).
\]

Instead of:

\begin{verbatim}
vector<vector<Gaussian>>
\end{verbatim}

GPU sees two flat buffers.

Much better for memory access.

\begin{center}\rule{0.5\linewidth}{0.5pt}\end{center}

\textbf{Maveb builds ranges after sorting}

It initializes every tile range to:

\[
(0xffffffff,0).
\]

Then walks sorted keys.

When tile ID changes, it closes previous range and opens next one.

Thus:

\protect\phantomsection\label{sxtu5b}%
\begin{verbatim}
sorted keys
→ contiguous same-tile runs
→ tile ranges
\end{verbatim}

No separate per-tile vector allocation exists.

\begin{center}\rule{0.5\linewidth}{0.5pt}\end{center}

After scan, Gaussian \(i\) has:

\[
destination=offset[i].
\]

Then loop over its tile rectangle:

\protect\phantomsection\label{4hi99e}%
\begin{verbatim}
for tileY:
    for tileX:
        values[destination] =
            gaussianIndex;

        keys[destination] =
            MakeSortKey(
                tile,
                depth);

        ++destination;
\end{verbatim}

Current Maveb does precisely this.

\begin{center}\rule{0.5\linewidth}{0.5pt}\end{center}

\textbf{Bounded capacity}

The total number of tile references can explode.

A huge Gaussian can touch thousands of tiles.

Millions of Gaussians can produce tens or hundreds of millions of references.

So Maveb gives the pipeline an explicit budget:

\[
maximumTileEntries.
\]

Default currently:

\[
4,194,304.
\]

Writes are guarded:

\begin{verbatim}
if (
    destination <
    maximumTileEntries)
{
    ...
}
\end{verbatim}

Overflow amount is counted.

This is critical.

Never let a projected Gaussian field turn into an out-of-bounds GPU write.

\begin{center}\rule{0.5\linewidth}{0.5pt}\end{center}

\textbf{Saturated work}

Maveb's scan/prefix path tracks:

\[
requested
\]

versus:

\[
stored=
\min(requested,capacity).
\]

It increments:

\protect\phantomsection\label{hhf4kg}%
\begin{verbatim}
overflowedEntries
\end{verbatim}

for excess references.

So failure mode becomes:

\[
\boxed{
\text{bounded visual degradation + telemetry}
}
\]

rather than:

\[
\boxed{
\text{memory corruption}.
}
\]

That is production engineering.

\begin{center}\rule{0.5\linewidth}{0.5pt}\end{center}

Each tile reference needs two ordering dimensions:

\[
tileID
\]

and:

\[
depth.
\]

Ultimately we want:

\protect\phantomsection\label{rgchv5}%
\begin{verbatim}
Tile 0:
    near → far

Tile 1:
    near → far

Tile 2:
    near → far
\end{verbatim}

So logically key is:

\[
\boxed{
(tileID,depth).
}
\]

Maveb stores:

\begin{verbatim}
uint2 key
\end{verbatim}

as:

\begin{verbatim}
key.x =
    as_type<uint>(depth);

key.y =
    tileID;
\end{verbatim}

At first glance this looks reversed.

But radix-pass ordering makes it correct.

\begin{center}\rule{0.5\linewidth}{0.5pt}\end{center}

\textbf{Why positive float bits can serve as depth key}

IEEE-754 positive finite floats preserve numerical order under unsigned bit interpretation.

If:

\[
0<a<b,
\]

then for ordinary positive finite values:

\[
bits(a)<bits(b).
\]

Maveb's camera depth is positive:

\[
Z>near>0.
\]

Therefore:

\begin{verbatim}
as_type<uint>(depth)
\end{verbatim}

can be sorted numerically as unsigned integer.

If negative depths were allowed, this trick would fail without sign-bit transformation.

But they have already been culled.

\begin{center}\rule{0.5\linewidth}{0.5pt}\end{center}

Suppose foreground A:

\[
C_A,\alpha_A
\]

and background B:

\[
C_B,\alpha_B.
\]

A over B gives:

\[
C=
\alpha_A C_A
+
(1-\alpha_A)
\alpha_B C_B.
\]

Reverse:

\[
C'=
\alpha_B C_B
+
(1-\alpha_B)
\alpha_A C_A.
\]

Generally:

\[
\boxed{
C\ne C'.
}
\]

Therefore alpha blending is not commutative.

Depth order matters.

\begin{center}\rule{0.5\linewidth}{0.5pt}\end{center}

\textbf{Concrete example}

Red A:

\[
C_A=(1,0,0),
\quad
\alpha_A=0.5.
\]

Blue B:

\[
C_B=(0,0,1),
\quad
\alpha_B=0.5.
\]

A in front:

\[
C=
0.5(1,0,0)
+
0.5\cdot0.5(0,0,1)
\]

\[
=
(0.5,0,0.25).
\]

B in front:

\[
C'=
(0.25,0,0.5).
\]

Very different.

Sorting is not optional.

\begin{center}\rule{0.5\linewidth}{0.5pt}\end{center}

One option:

Sort every visible Gaussian globally by depth:

\[
G_{near}\rightarrow G_{far}.
\]

Then each tile can reference Gaussians in that order.

This is conceptually simple.

But if we duplicate a Gaussian across tiles anyway, sorting expanded references can directly produce per-tile contiguous ranges.

Maveb sorts:

\[
(tile,depth)
\]

references.

That simultaneously groups by tile and orders by depth.

\begin{center}\rule{0.5\linewidth}{0.5pt}\end{center}

\textbf{Why per-tile ordering is enough}

A pixel only receives Gaussians whose footprints overlap its tile.

Therefore relative ordering of:

\[
G_A
\]

in tile 12 and:

\[
G_B
\]

in tile 98 is irrelevant.

They can never affect the same pixel.

So a global total ordering contains unnecessary ordering constraints.

The true requirement is:

\[
\boxed{
\text{all splats that may affect the same pixel must be correctly ordered}.
}
\]

Tile-local sort is a conservative way to ensure that.

\begin{center}\rule{0.5\linewidth}{0.5pt}\end{center}

Suppose two splats have equal depth:

\[
z_A=z_B.
\]

Their alpha blend order still potentially changes output.

If GPU sorting arbitrarily permutes equal keys each frame, you can get:

\begin{itemize}
\tightlist
\item
  flickering,
\item
  nondeterministic captures,
\item
  CPU/GPU mismatch,
\item
  difficult debugging.
\end{itemize}

So sort should be stable.

A stable sort preserves earlier order among equal keys.

\begin{center}\rule{0.5\linewidth}{0.5pt}\end{center}

\textbf{Maveb's stability}

References are generated deterministically:

Gaussian index ascending.

Within each Gaussian, tile Y then tile X ascending.

The radix scatter uses prefix ranks preserving input order within each bucket.

Thus equal sort keys preserve initial relative order.

That gives deterministic tie behavior.

This matters much more than it first appears.

\begin{center}\rule{0.5\linewidth}{0.5pt}\end{center}

Comparison sort costs roughly:

\[
O(M\log M).
\]

Keys are fixed-width integers.

Radix sort can achieve:

\[
O(kM)
\]

where \(k\) is number of digit passes.

\begin{center}\rule{0.5\linewidth}{0.5pt}\end{center}

\textbf{LSD radix sort}

Least-significant-digit radix sort processes low-order digit first.

For a multi-component composite key:

\[
(tile,depth),
\]

if stable sorting is used, sort \textbf{secondary} key first, then primary key.

We want tile primary, depth secondary.

So:

\[
\boxed{
\text{sort depth first}
\rightarrow
\text{sort tile second}.
}
\]

Because tile sort is stable, depth order inside each tile is preserved.

That is exactly Maveb's pass ordering.

\begin{center}\rule{0.5\linewidth}{0.5pt}\end{center}

\textbf{Maveb parallel radix}

Each \texttt{uint2} has:

\[
64\text{ bits total}.
\]

Default optimized path uses:

\[
4\text{-bit digits}.
\]

Each 32-bit component therefore needs:

\[
32/4=8
\]

passes.

Depth:

\[
8
\]

passes.

Tile:

\[
8
\]

passes.

Total:

\[
\boxed{16\text{ radix passes}.}
\]

The shader:

\begin{verbatim}
component =
    pass < 8
        ? key.x
        : key.y;
\end{verbatim}

is therefore deliberate.

\begin{center}\rule{0.5\linewidth}{0.5pt}\end{center}

\textbf{One radix pass}

For one digit with:

\[
16
\]

possible buckets, Maveb performs:

\protect\phantomsection\label{asjrz1}%
\begin{verbatim}
Histogram
↓
Global/group prefix offsets
↓
Stable scatter
\end{verbatim}

Histogram determines how many elements belong to each bucket.

Prefix determines output starting address.

Stable scatter places each key/value at destination.

This is again:

\[
\boxed{
\text{count}
\rightarrow
\text{scan}
\rightarrow
\text{scatter}.
}
\]

The exact same GPU pattern appears again.

\begin{center}\rule{0.5\linewidth}{0.5pt}\end{center}

\textbf{SIMD-group acceleration}

Maveb's histogram uses:

\begin{verbatim}
simd_sum()
\end{verbatim}

for bucket counts.

Scatter uses:

\begin{verbatim}
simd_prefix_exclusive_sum()
\end{verbatim}

for rank inside the SIMD group.

For one element belonging to bucket \(b\), local destination is:

\[
base_{group,b}
+
rank_{withinWorkgroup,b}.
\]

This avoids one atomic per sorted element.

\begin{center}\rule{0.5\linewidth}{0.5pt}\end{center}

\textbf{Why ping-pong buffers}

Sorting cannot safely scatter arbitrary elements in-place.

So Maveb owns:

\protect\phantomsection\label{t35rr1}%
\begin{verbatim}
Keys A
Keys B

Values A
Values B.
\end{verbatim}

Each pass:

\[
A\rightarrow B
\]

then:

\[
B\rightarrow A.
\]

After an even number of passes:

\[
16,
\]

result returns to A.

That is why range construction later reads:

\begin{verbatim}
keysA_
\end{verbatim}

and composite reads:

\begin{verbatim}
valuesA_.
\end{verbatim}

\begin{center}\rule{0.5\linewidth}{0.5pt}\end{center}

\textbf{Serial radix fallback}

Fallback uses:

\[
8\text{-bit digits}.
\]

Each 32-bit component:

\[
4
\]

passes.

Two components:

\[
8
\]

passes.

It is intentionally serial and deterministic.

Again:

\[
\text{fallback exists for correctness},
\]

not performance claims.

\begin{center}\rule{0.5\linewidth}{0.5pt}\end{center}

For pixel:

\[
p=
(x+0.5,y+0.5)
\]

and Gaussian center:

\[
\mu=(c_x,c_y),
\]

define:

\[
\boxed{
d=
p-\mu.
}
\]

So:

\[
d_x=
x+0.5-c_x,
\]

\[
d_y=
y+0.5-c_y.
\]

Maveb computes exactly this.

Pixel center convention:

\[
+0.5
\]

is important.

Do not evaluate at integer pixel corner unless the entire renderer is defined that way.

\begin{center}\rule{0.5\linewidth}{0.5pt}\end{center}

Conic:

\[
Q=
\Sigma^{-1}
=
\begin{bmatrix}
A&B\\
B&C
\end{bmatrix}.
\]

Mahalanobis squared distance:

\[
\boxed{
q=
d^TQd.
}
\]

Expanding:

\[
\boxed{
q=
Ad_x^2
+
2Bd_xd_y
+
Cd_y^2.
}
\]

Maveb:

\begin{verbatim}
distance =
    conic.x * dx * dx
  + 2 * conic.y * dx * dy
  + conic.z * dy * dy;
\end{verbatim}

\begin{center}\rule{0.5\linewidth}{0.5pt}\end{center}

Gaussian kernel:

\[
\boxed{
g(d)
=
e^{-\frac12q}.
}
\]

If:

\[
q=0,
\]

then:

\[
g=1.
\]

If:

\[
q=1,
\]

then:

\[
g\approx0.6065.
\]

If:

\[
q=4,
\]

then:

\[
g=e^{-2}\approx0.1353.
\]

If:

\[
q=9,
\]

then:

\[
g=e^{-4.5}\approx0.0111.
\]

Maveb rejects:

\[
q>9.
\]

This is the 3σ ellipse.

\begin{center}\rule{0.5\linewidth}{0.5pt}\end{center}

Current code:

\begin{verbatim}
exp(
    -0.5f *
    distance)
\end{verbatim}

is expensive relative to simple arithmetic, but only performed after:

\begin{itemize}
\tightlist
\item
  camera culling,
\item
  screen culling,
\item
  tile culling,
\item
  3σ conic rejection.
\end{itemize}

That ordering matters.

Do not compute exponent for obviously irrelevant splats.

\begin{center}\rule{0.5\linewidth}{0.5pt}\end{center}

Stored decoded opacity:

\[
\alpha_0.
\]

Pixel alpha:

\[
\boxed{
\alpha_i(p)
=
\alpha_{0,i}
e^{-q_i/2}.
}
\]

Maveb additionally caps:

\[
\boxed{
\alpha_i
=
\min(
0.99,
\alpha_{0,i}e^{-q/2}
).
}
\]

and rejects:

\[
\alpha_i<\frac1{255}.
\]

\begin{center}\rule{0.5\linewidth}{0.5pt}\end{center}

\textbf{Why cap at 0.99?}

Without cap, one Gaussian can produce:

\[
\alpha=1.
\]

Then transmittance becomes exactly:

\[
T=0.
\]

A cap preserves a small amount:

\[
T\ge0.01
\]

after one splat.

This avoids an individual primitive becoming mathematically perfectly opaque and improves stability around learned saturation.

It also mirrors standard 3DGS-style implementations where alpha is capped below one.

\begin{center}\rule{0.5\linewidth}{0.5pt}\end{center}

Start accumulated color:

\[
C_0=0.
\]

Start transmittance:

\[
T_0=1.
\]

For ordered Gaussian \(i\):

\[
\boxed{
w_i=
T_i\alpha_i.
}
\]

Then:

\[
\boxed{
C_{i+1}
=
C_i
+
w_i c_i.
}
\]

Update transmittance:

\[
\boxed{
T_{i+1}
=
T_i(1-\alpha_i).
}
\]

That is the central 3DGS compositing recurrence.

\begin{center}\rule{0.5\linewidth}{0.5pt}\end{center}

\textbf{Accumulated opacity formulation}

Instead of storing \(T\), Maveb stores accumulated opacity:

\[
A=1-T.
\]

Then:

\[
T=1-A.
\]

Contribution:

\[
\boxed{
w_i=
(1-A_i)\alpha_i.
}
\]

Update:

\[
A_{i+1}
=
A_i+w_i.
\]

Maveb code:

\begin{verbatim}
contribution =
    (1.0f - opacity)
    * alpha;

color +=
    contribution
    * gaussian.color;

opacity +=
    contribution;
\end{verbatim}

Exactly equivalent.

\begin{center}\rule{0.5\linewidth}{0.5pt}\end{center}

\textbf{Prove opacity recurrence}

Suppose:

\[
A_i=1-T_i.
\]

Then:

\[
A_{i+1}
=
A_i
+
(1-A_i)\alpha_i.
\]

Substitute:

\[
1-T_i
+
T_i\alpha_i
\]

\[
=
1-T_i(1-\alpha_i).
\]

Thus:

\[
T_{i+1}
=
T_i(1-\alpha_i).
\]

Correct.

\begin{center}\rule{0.5\linewidth}{0.5pt}\end{center}

After \(n\) splats:

\[
\boxed{
T_n
=
\prod_{i=0}^{n-1}
(1-\alpha_i).
}
\]

Therefore accumulated opacity:

\[
\boxed{
A_n
=
1-
\prod_i(1-\alpha_i).
}
\]

Contribution of splat \(i\):

\[
\boxed{
w_i
=
\alpha_i
\prod_{j<i}
(1-\alpha_j).
}
\]

This formula is worth memorizing.

It says:

\begin{quote}
A Gaussian contributes its own alpha multiplied by how much light remains visible through everything in front of it.
\end{quote}

\begin{center}\rule{0.5\linewidth}{0.5pt}\end{center}

\textbf{Background}

After all splats, background contribution is:

\[
\boxed{
C_{bg}^{contribution}
=
T_{final}C_{bg}.
}
\]

Thus final color:

\[
\boxed{
C=
\sum_i
T_i\alpha_i c_i
+
T_{final}C_{bg}.
}
\]

Maveb CPU oracle adds exactly:

\[
(1-A)C_{background}.
\]

The Metal Gaussian target stores its accumulated Gaussian color/alpha and later scene composition can determine how it participates in the wider render.

\begin{center}\rule{0.5\linewidth}{0.5pt}\end{center}

If:

\[
T
\]

becomes tiny, remaining splats cannot substantially change pixel.

Maveb terminates when:

\[
A>0.999.
\]

Equivalent:

\[
\boxed{
T<0.001.
}
\]

Remaining maximum contribution from bounded color \(C_{max}\):

\[
\le
0.001C_{max}.
\]

That gives an actual error interpretation.

\begin{center}\rule{0.5\linewidth}{0.5pt}\end{center}

\textbf{Early termination saves exactly the expensive tail}

Dense scenes often have many splats behind already-opaque foreground.

Without early out:

\protect\phantomsection\label{ss9s8z}%
\begin{verbatim}
evaluate
exp
blend
evaluate
exp
blend
...
\end{verbatim}

for potentially hundreds/thousands of hidden splats.

With front-to-back compositing, once \(T\) is tiny:

\[
\boxed{\text{stop}.}
\]

This is one reason front-to-back ordering is operationally valuable.

Maveb counts:

\protect\phantomsection\label{7ar1nv}%
\begin{verbatim}
earlyTerminations
\end{verbatim}

in GPU statistics.

That is useful performance evidence.

\begin{center}\rule{0.5\linewidth}{0.5pt}\end{center}

Traditional ``over'' can also be evaluated back-to-front.

Start with background:

\[
C=C_{bg}.
\]

For splats from far to near:

\[
\boxed{
C
\leftarrow
\alpha_i c_i
+
(1-\alpha_i)C.
}
\]

This yields the same ideal result as front-to-back when order is exact and arithmetic exact.

\begin{center}\rule{0.5\linewidth}{0.5pt}\end{center}

\textbf{Why front-to-back is preferable here}

Back-to-front cannot easily stop because near splats have not yet been processed.

A background splat that looks important may later become completely occluded.

Front-to-back knows remaining transmittance directly.

So:

\[
\boxed{
\text{front-to-back}
\rightarrow
\text{natural early termination}.
}
\]

For millions of semi-transparent splats, this matters.

\begin{center}\rule{0.5\linewidth}{0.5pt}\end{center}

Define premultiplied sample color:

\[
\tilde C_i
=
\alpha_i c_i.
\]

Then front-to-back:

\[
C
\leftarrow
C+
T\tilde C_i.
\]

and:

\[
T
\leftarrow
T(1-\alpha_i).
\]

Premultiplication cleanly separates:

\[
\text{color energy}
\]

from:

\[
\text{remaining transparency}.
\]

Maveb effectively computes premultiplied contribution on the fly:

\begin{verbatim}
contribution *
gaussian.color
\end{verbatim}

where:

\[
contribution=T\alpha.
\]

\begin{center}\rule{0.5\linewidth}{0.5pt}\end{center}

\textbf{Why straight-color interpolation causes edge problems}

If transparent colors are stored as straight RGB, fully transparent pixels can carry arbitrary RGB.

Filtering/interpolation may pull those hidden colors into visible borders.

Premultiplied representation satisfies:

\[
\alpha=0
\Rightarrow
\tilde C=0.
\]

This generally behaves better under filtering/compositing.

\begin{center}\rule{0.5\linewidth}{0.5pt}\end{center}

Naively:

\[
e^{-q/2}
\]

can underflow for huge:

\[
q.
\]

For float32, sufficiently negative exponents become subnormal then zero.

But Maveb does something better:

\begin{verbatim}
if (distance > 9)
    continue;

alpha =
    opacity *
    exp(-0.5 * distance);
\end{verbatim}

Therefore exponent argument always lies:

\[
[-4.5,0]
\]

for evaluated samples.

No meaningful exponential underflow can occur.

This is a perfect example of:

\[
\boxed{
\text{algorithmic culling improves numerical robustness as well as performance}.
}
\]

\begin{center}\rule{0.5\linewidth}{0.5pt}\end{center}

\textbf{Transmittance underflow}

Another issue is:

\[
T=
\prod_i(1-\alpha_i).
\]

After many splats, \(T\) may become extremely small.

But Maveb stops at:

\[
T<10^{-3}.
\]

So it never multiplies through thousands of values toward true floating-point underflow.

Again:

\[
\boxed{
\text{early termination prevents both wasted work and pathological numerics}.
}
\]

\begin{center}\rule{0.5\linewidth}{0.5pt}\end{center}

\subsection{THE M148 DATAFLOW IN CURRENT MAVEB}\label{04-visual-computing-ii-ch039-the-m148-dataflow-in-current-maveb}

\hypertarget{04-visual-computing-ii-ch039-the-m148-dataflow-in-current-maveb}{}

The current Metal pipeline owns:

\protect\phantomsection\label{f7zvxr}%
\begin{verbatim}
Projected Gaussians
Tile Counts
Tile Offsets
Scan Block Sums
Radix Histograms
Radix Group Offsets
Indirect Dispatch
Keys A/B
Values A/B
Tile Ranges
Counters
\end{verbatim}

and executes:

\protect\phantomsection\label{1ly8ek}%
\begin{verbatim}
Counter Clear
↓
Gaussian Projection
↓
Block Scan
↓
Block Prefix
↓
Block Offset Propagation
↓
Key Generation
↓
16×
    Radix Histogram
    Radix Group Prefix
    Stable Scatter
↓
Range Initialization
↓
Range Construction
↓
Tile Composite
\end{verbatim}

This is not a conceptual example.

It is the inspected Maveb production path.

\begin{center}\rule{0.5\linewidth}{0.5pt}\end{center}

\subsection{CPU ORACLE, METAL EXECUTION AND QUANTITATIVE AGREEMENT}\label{04-visual-computing-ii-ch039-cpu-oracle-metal-execution-and-quantitative-agreement}

\hypertarget{04-visual-computing-ii-ch039-cpu-oracle-metal-execution-and-quantitative-agreement}{}

Maveb's CPU oracle first projects all valid Gaussians.

Then:

\begin{verbatim}
std::ranges::stable_sort(
    projected,
    {},
    &Projected::depth);
\end{verbatim}

Since reference camera uses:

\[
+Z
\]

forward and smaller positive \(Z\) is closer:

\[
\boxed{
\text{ascending depth}
=
\text{front-to-back}.
}
\]

Stable sort preserves source order for equal depth.

This is deliberately compatible with deterministic GPU ordering.

\begin{center}\rule{0.5\linewidth}{0.5pt}\end{center}

\textbf{Why a CPU oracle should be simple}

The CPU reference implementation should optimize primarily for:

\[
\boxed{
\text{clarity + determinism + correctness}.
}
\]

Not maximum performance.

If CPU and GPU share the same clever algorithm and same bug, comparison proves nothing.

A useful oracle should be independent enough that implementation mistakes are unlikely to correlate.

Maveb's CPU rasterizer does not run the GPU's radix pipeline.

It globally stable-sorts projected records, then rasterizes them.

Good.

\begin{center}\rule{0.5\linewidth}{0.5pt}\end{center}

\textbf{Yet CPU still avoids pathological \(P\times N\)}

Maveb partitions the image into horizontal:

\[
16\text{-row bands}.
\]

For each projected Gaussian it computes pixel AABB and inserts a \texttt{BandEntry} into every intersected horizontal band.

Then each band processes only relevant Gaussians.

This is not the exact same 16×16 tile algorithm as GPU.

That difference is useful.

\begin{center}\rule{0.5\linewidth}{0.5pt}\end{center}

\textbf{Why deterministic threading survives}

Global projected list is already sorted.

Band lists receive Gaussians in that sorted order.

Every band owns disjoint rows.

So threads never concurrently update the same pixel.

Thus:

\[
\boxed{
\text{thread scheduling does not affect pixel accumulation order}.
}
\]

This gives deterministic multi-threaded CPU output.

That is an excellent oracle design.

\begin{center}\rule{0.5\linewidth}{0.5pt}\end{center}

A CPU implementation could use the same conceptual tile ranges:

\protect\phantomsection\label{tfx8be}%
\begin{verbatim}
for each tile:
{
    const auto [begin,end] =
        ranges[tile];

    for each pixel in tile:
    {
        float T = 1.0f;
        Vec3 color{};

        for (
            reference = begin;
            reference < end;
            ++reference)
        {
            evaluate splat;

            if (outside ellipse)
                continue;

            ...

            if (T < threshold)
                break;
        }
    }
}
\end{verbatim}

This has excellent cache locality because one tile's Gaussian list is shared by its pixels.

Maveb's oracle currently uses horizontal bands instead because it is simpler and remains independent from GPU scheduling machinery.

That independence is valuable for validation.

\begin{center}\rule{0.5\linewidth}{0.5pt}\end{center}

Current shared CPU/MSL ABI:

\protect\phantomsection\label{zxyqtm}%
\begin{verbatim}
struct AetherGaussianGpu
{
    float4 positionOpacity;
    float4 logScaleRestCount;
    float4 rotation;
    float4 dc;
    float4 shRest[12];
};
\end{verbatim}

Size:

\[
\boxed{256\text{ bytes}.}
\]

Why 12 \texttt{float4}s?

There are:

\[
45
\]

possible higher SH floats.

Twelve float4s provide:

\[
48
\]

float slots.

Three padding positions remain.

This keeps 16-byte alignment straightforward.

\begin{center}\rule{0.5\linewidth}{0.5pt}\end{center}

\textbf{Packed semantic fields}

\texttt{positionOpacity}:

\[
(x,y,z,opacityLogit).
\]

\texttt{logScaleRestCount}:

\[
(\ell_x,\ell_y,\ell_z,restCount).
\]

\texttt{rotation}:

\[
(w,x,y,z).
\]

\texttt{dc}:

\[
(R_{DC},G_{DC},B_{DC},padding).
\]

\texttt{shRest{[}12{]}} stores the 45 higher coefficients.

This is a GPU-friendly fixed-size record.

\begin{center}\rule{0.5\linewidth}{0.5pt}\end{center}

Projection:

\[
\boxed{
1\ GPU\ thread
\leftrightarrow
1\ Gaussian.
}
\]

Kernel receives:

\begin{verbatim}
uint index
[[thread_position_in_grid]]
\end{verbatim}

and exits if:

\[
index\ge GaussianCount.
\]

This is embarrassingly parallel:

Each Gaussian independently computes:

\begin{itemize}
\tightlist
\item
  camera-space center,
\item
  covariance,
\item
  conic,
\item
  radius,
\item
  SH color,
\item
  tile bounds,
\item
  overlap count.
\end{itemize}

Only shared counter:

\[
visibleGaussians
\]

needs atomic increment.

\begin{center}\rule{0.5\linewidth}{0.5pt}\end{center}

\textbf{Projection threadgroup width}

Maveb's generic \texttt{dispatch1D()} chooses:

\[
\min(
pipeline.maxTotalThreadsPerThreadgroup(),
256
).
\]

This is a conservative implementation.

A more performance-tuned path might consider:

\[
threadExecutionWidth
\]

and occupancy, but 256 is a conventional GPU workgroup size.

\begin{center}\rule{0.5\linewidth}{0.5pt}\end{center}

\textbf{Scan mapping}

Scan:

\[
256\text{ threads/workgroup}.
\]

One element:

\[
\leftrightarrow
\]

one Gaussian count.

Threadgroup shared scratch:

\[
256\times4
=
1024\text{ bytes}.
\]

Very cheap.

\begin{center}\rule{0.5\linewidth}{0.5pt}\end{center}

\textbf{Radix mapping}

Radix histogram/scatter:

\[
256
\]

threads/group.

SIMD-group intrinsics compute local bucket rank/count.

This is exactly the SIMT material from Phase 22 applied to a real engine workload.

\begin{center}\rule{0.5\linewidth}{0.5pt}\end{center}

\textbf{Composite mapping}

Current Metal composite dispatch:

\begin{verbatim}
dispatchThreads(
    {width,height,1},
    {8,8,1});
\end{verbatim}

Therefore:

\[
\boxed{
1\ GPU\ thread
\leftrightarrow
1\ pixel.
}
\]

Each workgroup contains:

\[
8\times8=64
\]

threads.

\begin{center}\rule{0.5\linewidth}{0.5pt}\end{center}

\textbf{Relationship between 8×8 compute groups and 16×16 Gaussian tiles}

Gaussian scheduling tile:

\[
16\times16.
\]

Compute workgroup:

\[
8\times8.
\]

Therefore one Gaussian tile covers:

\[
2\times2
\]

compute workgroups.

All four workgroups independently calculate the same tile ID and read the same:

\[
range[tile].
\]

That is correct.

The scheduling tile size does not have to equal compute threadgroup dimensions.

They solve different problems.

\begin{center}\rule{0.5\linewidth}{0.5pt}\end{center}

Current projected record:

\protect\phantomsection\label{2lb4al}%
\begin{verbatim}
struct AetherProjectedGaussian
{
    float4 centerDepthRadius;
    float4 conicOpacity;
    float4 color;

    uint4 tileBounds;
    uint4 sourceCountValid;
};
\end{verbatim}

Size:

\[
\boxed{80\text{ bytes}.}
\]

\begin{center}\rule{0.5\linewidth}{0.5pt}\end{center}

\textbf{Field semantics}

\[
centerDepthRadius
=
(c_x,c_y,z,r).
\]

\[
conicOpacity
=
(A,B,C,\alpha_0).
\]

\[
color
=
(R,G,B,1).
\]

\texttt{tileBounds} contains:

\[
(minX,minY,maxX,maxY)
\]

in tile coordinates.

\texttt{sourceCountValid} currently stores:

\[
(sourceIndex,\ overlapCount,\ valid,\ SHdegree).
\]

This is a well-designed projection product because composite kernel no longer needs:

\begin{itemize}
\tightlist
\item
  world-space covariance,
\item
  quaternion,
\item
  log scale,
\item
  SH coefficients,
\item
  camera projection math.
\end{itemize}

All that expensive work was amortized once per Gaussian.

\begin{center}\rule{0.5\linewidth}{0.5pt}\end{center}

\textbf{Projection is a representation-lowering stage}

This resembles compiler lowering.

Rich representation:

\protect\phantomsection\label{jyhwui}%
\begin{verbatim}
3D mean
quaternion
scale
SH coefficients
opacity logit
\end{verbatim}

becomes:

\protect\phantomsection\label{z54oi2}%
\begin{verbatim}
2D center
depth
conic
opacity
color
tile bounds
\end{verbatim}

The second form is exactly what rasterization needs.

This is why GPU pipelines benefit from explicit intermediate representations.

\begin{center}\rule{0.5\linewidth}{0.5pt}\end{center}

There are at least two broad ways to splat.

Raster approach:

Create one quad per projected Gaussian.

Rasterize quad.

Fragment shader evaluates Gaussian.

Use blending.

Compute approach:

Tile/bin Gaussians.

Dispatch one thread per pixel.

Explicitly loop ordered splats and perform compositing.

Maveb currently chooses compute.

\begin{center}\rule{0.5\linewidth}{0.5pt}\end{center}

\textbf{Raster splat advantages}

Hardware rasterization gives:

\begin{itemize}
\tightlist
\item
  coverage generation,
\item
  primitive setup,
\item
  blend units,
\item
  easy quad projection.
\end{itemize}

A simple Gaussian renderer can be built quickly this way.

Vulkax even contains projected-raster paths using generated quad geometry.

\begin{center}\rule{0.5\linewidth}{0.5pt}\end{center}

\textbf{Raster disadvantages}

Correct ordering for millions of translucent splats remains difficult.

Blend hardware cannot magically solve arbitrary per-pixel ordering.

A global splat order may work approximately but has problems with overlapping anisotropic geometry and per-tile ordering.

You may also pay:

\begin{itemize}
\tightlist
\item
  huge quad overdraw,
\item
  fragment invocation outside true ellipse,
\item
  per-splat draw/instance management.
\end{itemize}

\begin{center}\rule{0.5\linewidth}{0.5pt}\end{center}

\textbf{Compute advantages}

Compute lets renderer explicitly control:

\[
\text{tile list}
\]

\[
\text{ordering}
\]

\[
\text{early termination}
\]

\[
\text{debug IDs/depth}
\]

\[
\text{overflow handling}.
\]

This makes it an excellent fit for serious 3DGS rasterization.

\begin{center}\rule{0.5\linewidth}{0.5pt}\end{center}

\textbf{Compute disadvantage}

You implement much more yourself:

\begin{itemize}
\tightlist
\item
  scan,
\item
  sort,
\item
  ranges,
\item
  pixel loops,
\item
  synchronization.
\end{itemize}

Which is precisely why Maveb has 13 compute pipelines for this subsystem.

\begin{center}\rule{0.5\linewidth}{0.5pt}\end{center}

Dependencies:

\protect\phantomsection\label{sd09sx}%
\begin{verbatim}
Projection writes:
    projected
    counts

Scan reads counts
writes offsets

Key generation reads:
    projected
    offsets
writes:
    keys
    values

Radix reads/writes:
    keys
    values

Range construction reads keys
writes ranges

Composite reads:
    projected
    sorted values
    ranges

Composite writes:
    color/depth/ID textures
\end{verbatim}

These are real RAW dependencies.

\begin{center}\rule{0.5\linewidth}{0.5pt}\end{center}

\textbf{Current Metal synchronization strategy}

Maveb ends the compute encoder after each stage.

Then creates the next encoder on the \textbf{same command buffer}.

The relevant buffers use normal tracked Metal hazard semantics; the pipeline is not using an explicit untracked-memory scheme here.

Thus Metal orders and tracks producer→consumer dependencies across encoders.

There is no long sequence of Vulkan-style manual barriers in \texttt{GaussianPipeline::encode()}.

That does \textbf{not} mean synchronization is absent.

It means:

\[
\boxed{
\text{dependency is expressed through ordered encoders + Metal resource hazard tracking}.
}
\]

Exactly the Metal-vs-Vulkan distinction from Pass 13.

\begin{center}\rule{0.5\linewidth}{0.5pt}\end{center}

\textbf{Why separate encoders are convenient here}

Each stage has a clear boundary:

\protect\phantomsection\label{9acmg5}%
\begin{verbatim}
Projection encoder ends
↓
Scan encoder starts
\end{verbatim}

so Metal can understand that a producer phase precedes consumer phase.

Within one compute encoder, if dispatches read data written by previous dispatch without sufficient implicit scope, memory-barrier semantics may be needed.

Maveb's design avoids ambiguity by giving major stages separate encoders.

\begin{center}\rule{0.5\linewidth}{0.5pt}\end{center}

A serious GPU renderer needs an oracle.

Test conditions should use identical:

\[
\text{Gaussian asset},
\]

\[
\text{camera intrinsics/extrinsics},
\]

\[
\text{viewport},
\]

\[
\text{background},
\]

\[
\text{SH convention},
\]

\[
\text{alpha thresholds}.
\]

Then render:

\[
I_{CPU}
\]

and:

\[
I_{GPU}.
\]

Compare in:

\[
\boxed{\text{linear color space}.}
\]

Do not compare after arbitrary tone mapping if testing the Gaussian integrator itself.

\begin{center}\rule{0.5\linewidth}{0.5pt}\end{center}

\textbf{Maveb's current correctness philosophy}

The repository's accepted Metal Gaussian ADR says the standard renderer includes:

\begin{itemize}
\tightlist
\item
  projection,
\item
  anisotropic covariance,
\item
  culling,
\item
  overlap counts,
\item
  exclusive offsets,
\item
  key generation,
\item
  stable ordering,
\item
  tile ranges,
\item
  front-to-back composition,
\item
  depth,
\item
  IDs,
\item
  early termination,
\item
  bounded overflow.
\end{itemize}

It also specifies a Metal test spanning more than one scan/radix block with shuffled depths and varying colors, compared against the CPU oracle.

That is exactly the kind of test needed.

I did \textbf{not} find a current repository implementation named specifically around all four curriculum metrics---RMSE, PSNR, maximum error and changed-pixel fraction---in the inspected Gaussian test paths, so the formulas below are the validation methodology required by the curriculum rather than a claim that Maveb already exposes each metric as a production API.

\begin{center}\rule{0.5\linewidth}{0.5pt}\end{center}

For two images with:

\[
N
\]

scalar channel samples:

\[
e_i=
G_i-R_i.
\]

MSE:

\[
\boxed{
MSE=
\frac1N
\sum_{i=1}^{N}
e_i^2.
}
\]

RMSE:

\[
\boxed{
RMSE=
\sqrt{
\frac1N
\sum_i
(G_i-R_i)^2
}.
}
\]

Units are same as image values.

If image is linear normalized RGB, RMSE of:

\[
10^{-3}
\]

means average magnitude of numerical difference is small.

But exact acceptable threshold depends on scene, precision and expected backend differences.

Do not invent universal threshold.

\begin{center}\rule{0.5\linewidth}{0.5pt}\end{center}

\textbf{RGB RMSE}

For \(P\) pixels:

\[
\boxed{
RMSE_{RGB}
=
\sqrt{
\frac{
\sum_{p}
[
(\Delta R_p)^2+
(\Delta G_p)^2+
(\Delta B_p)^2
]
}{
3P
}
}.
}
\]

You may separately compute alpha RMSE.

Depth should have its own metric.

IDs should generally be discrete agreement, not RMSE.

\begin{center}\rule{0.5\linewidth}{0.5pt}\end{center}

For range:

\[
[0,MAX],
\]

PSNR:

\[
\boxed{
PSNR
=
10\log_{10}
\frac{MAX^2}{MSE}.
}
\]

If:

\[
MAX=1,
\]

then:

\[
PSNR
=
-10\log_{10}(MSE).
\]

If MSE:

\[
10^{-4},
\]

then:

\[
PSNR=40\text{ dB}.
\]

If:

\[
10^{-6},
\]

then:

\[
60\text{ dB}.
\]

Higher is better.

\begin{center}\rule{0.5\linewidth}{0.5pt}\end{center}

\textbf{PSNR caveat with HDR}

If HDR values exceed:

\[
1,
\]

what is:

\[
MAX?
\]

You must define it.

Possible approaches:

\begin{itemize}
\tightlist
\item
  fixed physical/reference range,
\item
  measured reference peak,
\item
  tone-mapped LDR metric.
\end{itemize}

Do not blindly use:

\[
MAX=1
\]

for unclamped HDR.

For backend numerical validation, RMSE plus max error is often less ambiguous.

\begin{center}\rule{0.5\linewidth}{0.5pt}\end{center}

Define:

\[
\boxed{
E_{\max}
=
\max_{p,c}
|G_{p,c}-R_{p,c}|.
}
\]

RMSE can hide one catastrophic pixel.

Example:

One pixel error:

\[
1
\]

among a million perfect pixels gives small RMSE.

But max error reveals:

\[
1.
\]

So both matter.

\begin{center}\rule{0.5\linewidth}{0.5pt}\end{center}

\textbf{Location of max error}

Do not store only scalar max.

Store:

\protect\phantomsection\label{vgsyyh}%
\begin{verbatim}
struct MaxError
{
    float Error;
    uint32_t X;
    uint32_t Y;
    uint32_t Channel;
};
\end{verbatim}

Then debugger can inspect:

\begin{itemize}
\tightlist
\item
  Gaussian ID,
\item
  depth,
\item
  tile,
\item
  sort rank,
\item
  conic.
\end{itemize}

That turns a metric into a diagnostic.

\begin{center}\rule{0.5\linewidth}{0.5pt}\end{center}

Choose threshold:

\[
\epsilon.
\]

Define changed indicator:

\[
I_p=
\begin{cases}
1,&
\max_c|\Delta_{p,c}|>\epsilon\\
0,&
\text{otherwise}.
\end{cases}
\]

Then:

\[
\boxed{
F_{changed}
=
\frac1P
\sum_p I_p.
}
\]

This distinguishes:

\begin{quote}
tiny errors everywhere
\end{quote}

from:

\begin{quote}
large/localized errors in one region.
\end{quote}

\begin{center}\rule{0.5\linewidth}{0.5pt}\end{center}

\textbf{Why RMSE + max + changed fraction is much more informative}

Case A:

Every pixel differs by:

\[
10^{-4}.
\]

Low RMSE.

Low max.

Changed fraction maybe 100\% for tiny epsilon.

Likely systematic floating-point variation.

Case B:

Only 1\% pixels wrong by:

\[
0.5.
\]

RMSE might still look moderate.

Max huge.

Changed fraction 1\%.

Likely culling/boundary/sorting issue.

Metrics together reveal error structure.

\begin{center}\rule{0.5\linewidth}{0.5pt}\end{center}

\textbf{Validate semantic buffers separately}

Maveb already produces:

\[
color,
depth,
ID.
\]

Its debug modes expose:

\protect\phantomsection\label{zp6hzo}%
\begin{verbatim}
appearance
logarithmic depth
source ID
tile occupancy
accumulated opacity
SH bands
sort rank
\end{verbatim}

That is excellent for fault isolation.

Color mismatch with matching depth/ID suggests:

\[
SH/alpha.
\]

ID mismatch suggests:

\[
ordering/culling.
\]

Tile occupancy mismatch suggests:

\[
projection/bounds/scheduler.
\]

Sort-rank mismatch suggests:

\[
key/radix/ranges.
\]

\begin{center}\rule{0.5\linewidth}{0.5pt}\end{center}

\textbf{ID validation}

Away from equal-depth/raster-boundary ambiguities, dominant source ID should often match exactly:

\[
\boxed{
ID_{CPU}(p)=ID_{GPU}(p).
}
\]

A color tolerance cannot excuse a completely different dominant Gaussian unless there is a documented near-tie.

This is an extremely strong validation signal.

\begin{center}\rule{0.5\linewidth}{0.5pt}\end{center}

\textbf{Depth validation}

Maveb stores depth of first accepted contributing splat:

CPU:

\begin{verbatim}
if depth == infinity:
    depth =
        gaussian.depth;
\end{verbatim}

GPU:

\begin{verbatim}
if (!isfinite(depth))
    depth =
        gaussian.centerDepthRadius.z;
\end{verbatim}

So CPU/GPU should agree closely if sort and alpha thresholds agree.

If they do not:

\[
\boxed{
\text{suspect order before shading}.
}
\]

\begin{center}\rule{0.5\linewidth}{0.5pt}\end{center}

\section{Give captured Gaussians stable world identity}\label{04-visual-computing-ii-ch039-give-captured-gaussians-stable-world-identity}

\hypertarget{04-visual-computing-ii-ch039-give-captured-gaussians-stable-world-identity}{}

A splat renderer answers an image question: which color and opacity does this primitive contribute at a pixel? Editing a captured scene asks a different question: which persistent world element does this primitive represent, and which physical variables are allowed to move it? These representations overlap, but they are not interchangeable. The original 3D Gaussian Splatting work optimizes anisotropic primitives for view synthesis and rendering; the entity and motion contract below is an explicit engine extension, not a property guaranteed by a trained PLY file. \href{https://arxiv.org/abs/2308.04079}{The 3D Gaussian Splatting paper}

Keep identity separate from storage order. Let each Gaussian carry an immutable identifier \(g_i=(n_i,\ell_i)\), where \(n_i\) is a namespace or asset identifier and \(\ell_i\) is a local identifier. The current vector index is only a temporary address. Sorting, compaction, filtering, or cloning may change that address; none of those operations may silently change the identity of an existing record. A lookup table maps \(g_i\) to the current index and is rebuilt after a reorder. A newly cloned primitive receives a new identifier and records its parent ID; it does not reuse the parent's ID.

An editable scene needs a second relation from Gaussian IDs to semantic entities and physical degrees of freedom. Represent it as a versioned association table. For Gaussian \(i\) and entity or node \(j\), let \(w_{ij}\) be the authored influence weight. For a Gaussian fully assigned to one rigid object, one weight is \(1\). For a soft boundary, several weights may be nonzero. Within one declared support set, require

\[
w_{ij}\ge 0,\qquad \sum_j w_{ij}=1.
\]

Normalization is necessary to reproduce a common translation. Reproducing a complete affine map also requires the rest-center reconstruction below, or an explicitly retained local offset. It does not prove that the association is correct: confidence, source, and an unassigned state must remain explicit. Do not turn a low-confidence semantic guess into a collider or a destructive edit automatically.

First consider weights with \textbf{linear precision}: the support positions reproduce the actual learned center, rather than merely averaging some nearby nodes. With associated rest positions \(x_j\), this means

\[
\mu_i=\sum_jw_{ij}x_j.
\]

If every support point receives the same affine map \(x'_j=Ax_j+b\), then

\[
\mu'_i=\sum_jw_{ij}(Ax_j+b)
=A\sum_jw_{ij}x_j+b\sum_jw_{ij}
=A\mu_i+b.
\]

The last equality uses \(\sum_jw_{ij}=1\). This is the affine-reproduction property: the weighted correspondence follows any shared translation, rotation, scale, or shear exactly. If the weights do not sum to one, even a common translation is scaled incorrectly.

A rigid attachment is not the same as replacing every center by the body's center of mass. For body center \(x_j\) and rotation \(R_j\), retain each splat's own local offset \(r_i=R_j^T(\mu_i-x_j)\). Its new center is \(x'_j+R'_jr_i\). A one-weight row says which body owns it; that row alone does not encode its distinct location.

For general support weights that do not reconstruct the learned center, retain \(d_i=\mu_i-\sum_jw_{ij}x_j\). Given a declared local affine map \(A_i\), use \(\mu'_i=\sum_jw_{ij}x'_j+A_id_i\). Under a common edit \(x'_j=Ax_j+b\) and \(A_i=A\), this still yields \(A\mu_i+b\). Dropping \(d_i\) would cause an immediate jump even before any edit. An empty or unreliable association remains explicitly unassigned; it should not be repaired by inventing uniform weights.

A Gaussian's covariance describes its spatial footprint. Under the same local linear map \(A\), transform it by

\[
\Sigma'_i=A\Sigma_iA^T.
\]

To derive this, write a centered sample as \(\delta=x-\mu\) with covariance \(E[\delta\delta^T]=\Sigma\). The transformed centered sample is \(\delta'=A\delta\); therefore

\[
E[\delta'\delta'^T]=E[A\delta\delta^TA^T]=A\Sigma A^T.
\]

Translation cancels from covariance. A nonsingular \(A\) preserves positive definiteness. If \(A\) collapses an axis, \(\Sigma'\) becomes singular and cannot be inverted to form a conic. A production edit must then reject the update or apply a separately declared minimum-scale policy; silently inverting a singular covariance is not a valid repair. For a pure translation, \(A=I\), so the mean moves and the covariance stays unchanged. For a rigid rotation \(R\), the covariance rotates as \(R\Sigma R^T\). For a uniform scale \(s\), it becomes \(s^2\Sigma\).

For a deforming support, estimate a local map only when the neighborhood contains enough independent motion. With centered rest offsets \(r_j=x_j-\mu\) and changed offsets \(r'_j=x'_j-\mu'\), a weighted least-squares map minimizes \(\sum_jw_{ij}\|r'_j-Ar_j\|^2\). If the weighted rest covariance \(M=\sum_jw_{ij}r_jr_j^T\) is nonsingular, the solution is

\[
A_i=\left(\sum_jw_{ij}r'_jr_j^T\right)M^{-1}.
\]

This formula follows by setting the matrix derivative of the squared residual to zero. Solve the resulting linear system with a rank-aware method rather than forming the inverse explicitly in numerical code. It supplies a local affine approximation over this primitive; spatially varying influences or a strongly curved deformation are not represented exactly by one matrix. If \(M\) is rank deficient, the neighborhood does not determine a full 3D affine map; use a lower-dimensional model, add a documented regularizer, or reject the deformation. The equation defines a local approximation, not a guarantee that arbitrary cloth folds or topology changes are represented faithfully.

\subsection{Worked edit: a supported Gaussian under stretch and translation}\label{04-visual-computing-ii-ch039-worked-edit-a-supported-gaussian-under-stretch-and-translation}

\hypertarget{04-visual-computing-ii-ch039-worked-edit-a-supported-gaussian-under-stretch-and-translation}{}

Two support nodes have rest positions \(x_1=(0,0)^T\) and \(x_2=(4,2)^T\). A Gaussian uses \(w_1=w_2=0.5\), so its mean is

\[
\mu=0.5(0,0)+0.5(4,2)=(2,1).
\]

Apply \(A=\begin{bmatrix}2&0\\0&1\end{bmatrix}\) and \(b=(1,-1)^T\). The nodes become \(x'_1=(1,-1)\) and \(x'_2=(9,1)\). The weighted mean is \((5,0)\), which agrees with \(A\mu+b=(5,0)\). Now let the original covariance be \(\Sigma=\begin{bmatrix}4&1\\1&2\end{bmatrix}\). Then

\[
\Sigma'=A\Sigma A^T=
\begin{bmatrix}2&0\\0&1\end{bmatrix}
\begin{bmatrix}4&1\\1&2\end{bmatrix}
\begin{bmatrix}2&0\\0&1\end{bmatrix}
=
\begin{bmatrix}16&2\\2&2\end{bmatrix}.
\]

Here \(A\) is supplied as the edit operator; two support nodes alone would not determine an unrestricted two-dimensional affine fit. The determinant changes from \(7\) to \(28\), matching \(\det(A)^2\det(\Sigma)=4\cdot7\). Translation changed the center but contributed nothing to the covariance. A student can check the result by multiplying the matrices and confirming \(\det(\Sigma')>0\).

\section{Couple edits in both directions and keep their proof}\label{04-visual-computing-ii-ch039-couple-edits-in-both-directions-and-keep-their-proof}

\hypertarget{04-visual-computing-ii-ch039-couple-edits-in-both-directions-and-keep-their-proof}{}

An appearance edit and a physical edit are separate operations with an explicit bridge between them. In one direction, a semantic selection resolves stable entity IDs to stable Gaussian IDs, then maps those Gaussian IDs through the versioned support table to affected physical DOFs. In the other direction, a committed physical state computes the local motion map for each supported Gaussian, updates its mean and covariance, and rebuilds projected bounds and tile references before rendering. A translation need not change an object-attached appearance function in this model. For spherical-harmonic (SH) appearance attached to a frame \(F\), evaluate its coefficients using local viewing direction \(F^Td_{world}\), and rotate \(F\) with the object. A renderer that instead stores world-direction coefficients must explicitly choose whether to preserve that world-space function or rotate its basis to follow the object. Deformation and relighting do not have an automatic appearance solution. Keep the appearance frame separate from the covariance eigenbasis: repeated eigenvalues and eigenvector sign choices can change that basis without changing the Gaussian shape.

To store a general deformed covariance back as rotation and scales, diagonalize the symmetric \(\Sigma'=Q\Lambda Q^T\), require positive eigenvalues, and use scales \(\sqrt{\lambda_j}\). Choose eigenvector signs so \(\det Q=+1\) before converting to a quaternion. At repeated eigenvalues, orientation is not unique; preserve continuity or retain the full covariance rather than interpreting an arbitrary eigenvector flip as physical rotation. Covariance push-forward changes shape, not a conserved physical mass. If opacity is an appearance peak it may be held fixed by policy; preserving an extinction or density integral needs a separately defined field and Jacobian rule. Do not silently assume either policy from \(A\Sigma A^T\).

A useful edit transaction has a before state, a proposed state, and a commit decision. Store the edit ID, parent scene revision, selected stable IDs, association-table revision, operator and parameters, before/after state hashes, affected Gaussian IDs, and validation results. Write the candidate state to a separate revision. Validate finite means, normalized orientation, positive-definite covariance, valid support weights, and stable ID uniqueness before publication. Render the candidate from the same cameras used for the before state. Compare the intended region and the declared unaffected region separately; include color, depth and dominant-ID checks where available. Commit the new revision only if all declared conditions pass. Otherwise retain the prior revision and the failed report. A provenance record makes a result repeatable; it does not make the physical interpretation true by itself.

\subsection{Reconstruction exercise with solution}\label{04-visual-computing-ii-ch039-reconstruction-exercise-with-solution}

\hypertarget{04-visual-computing-ii-ch039-reconstruction-exercise-with-solution}{}

A selected rigid cabinet has three associated Gaussians with stable IDs 51, 64 and 92. The cabinet is translated by \(b=(0,2,0)\) with \(A=I\). Their support weights each sum to one. For Gaussian 64, the old mean is \((3,1,0)\) and covariance is

\[
\Sigma=\begin{bmatrix}4&1&0\\1&2&0\\0&0&0.25\end{bmatrix}.
\]

The renderer's current vector indices for IDs 51, 64 and 92 are \((8,2,11)\), but a depth sort changes them to \((5,10,1)\). State the new mean, covariance, selected identity after sorting, and the minimum checks before committing.

\textbf{Solution.} Translation gives \(\mu' = \mu+b=(3,3,0)\). Since \(A=I\), \(\Sigma'=I\Sigma I^T=\Sigma\); translation does not change footprint. The selected identity remains ID 64, whose new transient index is 10. It is incorrect to preserve index 2 as the selection. Before commit, check IDs are unique, every affected support row remains nonnegative and sums to one, means are finite, each covariance is finite and positive definite, and the candidate render preserves declared unaffected regions within the chosen color/depth/ID tolerances. Save the new revision and its report only after those checks pass.

\subsection{Limits and exit check}\label{04-visual-computing-ii-ch039-limits-and-exit-check}

\hypertarget{04-visual-computing-ii-ch039-limits-and-exit-check}{}

A Gaussian is a learned appearance primitive, not automatically a surface point, rigid body, collider, or mass element. An association table is a model chosen for an application, and ambiguous or incomplete support should remain visible as uncertainty. The affine covariance rule preserves a Gaussian footprint under a locally affine map; it does not solve visibility, self-intersection, cloth topology, contact, or appearance changes caused by newly revealed surfaces. A render that matches the original cameras cannot prove an edit is physically correct from unseen views.

Before calling a captured-world edit complete, a learner should be able to: trace a semantic selection through stable IDs after reordering; show that normalized support weights reproduce a shared affine transform; derive and calculate \(A\Sigma A^T\); identify a rank-deficient local fit; rotate or preserve SH according to a stated frame; and produce a provenance record plus before/after tests for affected and unaffected regions. The original 3DGS paper establishes a view-synthesis representation and renderer, not these editor, physics, or publication contracts. Those are additional systems that must be designed and verified explicitly.

\section{Train the splats only after the forward renderer is understood}\label{04-visual-computing-ii-ch039-train-the-splats-only-after-the-forward-renderer-is-understood}

\hypertarget{04-visual-computing-ii-ch039-train-the-splats-only-after-the-forward-renderer-is-understood}{}

A projected, tiled, ordered, composited image is now a complete forward result, and the preceding CPU/GPU comparison gives us checks for color, depth, dominant ID, tile occupancy, and sorting. We can now ask the inverse question: how should the stored splats change so that this renderer matches calibrated photographs? This is where training belongs in the story: the forward map and its failure signals are already visible. The original 3D Gaussian Splatting paper starts from sparse points from camera calibration when available, optimizes their attributes, and interleaves optimization with density control. A rasterizer tutorial without this lifecycle explains only inference, not how the captured representation is fitted. (\href{https://arxiv.org/abs/2308.04079}{Kerbl et al., 2023 paper})

The author implementation used below is pinned to commit \href{https://github.com/graphdeco-inria/gaussian-splatting/commit/54c035f7834b564019656c3e3fcc3646292f727d}{\texttt{54c035f7834b564019656c3e3fcc3646292f727d}}, committed 30 October 2024. It is a precise historical source specimen, not a claim about every later fork or current 3DGS implementation. Its \href{https://github.com/graphdeco-inria/gaussian-splatting/blob/54c035f7834b564019656c3e3fcc3646292f727d/train.py\#L126-L186}{\texttt{train.py}} uses a weighted \(L_1\) and structural-similarity image loss, optionally adds depth regularization, differentiates the loss, gathers projected-mean gradient statistics, performs scheduled density control and opacity reset, then steps the optimizer. Its \href{https://github.com/graphdeco-inria/gaussian-splatting/blob/54c035f7834b564019656c3e3fcc3646292f727d/scene/gaussian_model.py\#L39-L47}{\texttt{GaussianModel}} maps stored scale through \(\exp\), opacity through \(\sigma\), and normalizes rotation quaternions. A robust importer or optimizer must still reject or explicitly repair a zero quaternion: normalization alone does not turn a zero vector into a valid rotation.

For calibrated camera \(k\), captured image \(I_k\), valid pixel set \(\Omega_k\), and renderer \(R_k(\Theta)\), a simplified scalar-channel objective is

\[
L(\Theta)=\sum_{k\in\mathcal B}\frac{1}{|\Omega_k|}\sum_{p\in\Omega_k}\rho\left(R_k(\Theta)_p-I_{k,p}\right),
\]

where \(\mathcal B\) is the sampled batch and \(\rho\) is the declared penalty. In the pinned training loop, the image term is a weighted \(L_1\) plus DSSIM loss; when reliable depth supervision is enabled, a weighted inverse-depth \(L_1\) term is added. The objective's exact weights and preprocessing are part of the experiment. Lower image error does not prove unique or physically useful geometry: several scenes can explain the same measurements, and a view-synthesis representation can exploit appearance to hide geometric error.

The reference stores several parameters in unconstrained form: mean \(\mu\in\mathbb R^3\); log-scale \(s_{log}\) decoded as \(s=\exp(s_{log})>0\); quaternion \(q\) normalized before it constructs a rotation; opacity logit \(a\) decoded as \(o=\sigma(a)\in(0,1)\) for finite real \(a\); and appearance coefficients such as spherical harmonics. These choices make ordinary optimizer steps compatible with positive scale and bounded opacity. The opacity derivative can still saturate near zero or one.

\section{Derive the image gradient through a Gaussian's footprint}\label{04-visual-computing-ii-ch039-derive-the-image-gradient-through-a-gaussians-footprint}

\hypertarget{04-visual-computing-ii-ch039-derive-the-image-gradient-through-a-gaussians-footprint}{}

At pixel \(x\), the Gaussian's spatial weight is \(G(x)=\exp(-q(x)/2)\), where \(q\) is its squared elliptical distance. Before implementation-specific clamping, its effective sample opacity is

\[
\alpha(x)=oG(x),\qquad o=\sigma(a).
\]

For a pixel where this unclamped branch is active and the footprint is held fixed,

\[
\frac{\partial\alpha}{\partial a}=G\,o(1-o).
\]

If the renderer clamps \(oG\) at a maximum, the derivative is zero on the saturated branch. The factor \(\alpha(1-\alpha)\) is only the special case \(G=1\); in general it is wrong because the Gaussian's footprint weight is separate from stored opacity.

For two ordered splats in one color channel over background \(c_b\), compositing gives

\[
C=\alpha_1c_1+(1-\alpha_1)\left[\alpha_2c_2+(1-\alpha_2)c_b\right].
\]

Holding order, centers, and footprint weights fixed, the color and front-opacity derivatives are

\[
\frac{\partial C}{\partial c_1}=\alpha_1,\qquad
\frac{\partial C}{\partial c_2}=(1-\alpha_1)\alpha_2,
\qquad
\frac{\partial C}{\partial\alpha_1}=c_1-\left[\alpha_2c_2+(1-\alpha_2)c_b\right].
\]

With squared error \(\ell=\tfrac12(C-y)^2\), the chain rule yields \(\partial\ell/\partial c_1=(C-y)\alpha_1\) and \(\partial\ell/\partial a_1=(C-y)(\partial C/\partial\alpha_1)G_1o_1(1-o_1)\) in the unclamped branch. Position and covariance gradients continue through the projected center and conic because changing them changes \(q(x)\) and thus \(G(x)\). Discrete tile membership, depth order, threshold crossings, and split/clone decisions are not ordinary smooth parameters; they form separate algorithmic boundaries.

A hand-checkable toy makes the chain rule concrete. Set both footprint weights to one, so the unclamped \(\alpha_i=o_i\); this is an explicit unit-footprint assumption, not the general pixel case. Use \(\alpha_1=o_1=0.5\), \(c_1=0.8\), \(\alpha_2=o_2=0.25\), \(c_2=0.2\), \(c_b=0\), target \(y=0.4\). Then \(C=0.425\), residual \(C-y=0.025\), and

\[
\frac{\partial\ell}{\partial c_1}=0.0125,\qquad
\frac{\partial\ell}{\partial\alpha_1}=0.01875,\qquad
\left.\frac{\partial\ell}{\partial a_1}\right|_{G_1=1}=0.0046875.
\]

A direct scalar SGD illustration with \(\eta=0.1\) gives \(c_1'=0.79875\). If alpha were optimized directly, \(\alpha_1'=0.498125\). With the stored logit, \(a_1'=a_1-0.00046875\); near \(a_1=0\), decoded opacity changes by about one quarter of that logit step because \(d\sigma/da=0.25\). A production run instead accumulates gradients from many pixels and camera views and uses adaptive optimizer state; this one-pixel trace checks signs and factors only.

\section{Density control changes model size between gradient steps}\label{04-visual-computing-ii-ch039-density-control-changes-model-size-between-gradient-steps}

\hypertarget{04-visual-computing-ii-ch039-density-control-changes-model-size-between-gradient-steps}{}

Continuous optimization cannot add a primitive when current support is inadequate. The pinned implementation accumulates the norm of the first two components of the projected-mean gradient for visible Gaussians and divides by a visibility count to estimate an average image-space signal (\href{https://github.com/graphdeco-inria/gaussian-splatting/blob/54c035f7834b564019656c3e3fcc3646292f727d/scene/gaussian_model.py\#L471-L472}{\texttt{add\_densification\_stats}}). At scheduled intervals it applies two capacity operations:

\begin{itemize}
\tightlist
\item
  \textbf{Clone:} when the accumulated mean-gradient magnitude reaches the threshold and the largest world-space scale is no greater than \texttt{percent\_dense\ *\ scene\_extent}, copy that small primitive's mean, appearance, opacity, scale, and rotation into a new primitive. The original remains. This adds local representation capacity without moving either copy.
\item
  \textbf{Split:} when the gradient reaches threshold but the largest scale exceeds that cutoff, replace the broad primitive with \(N\) children. In this pinned code \(N=2\) by default. It samples a zero-mean offset whose standard deviation is the parent's three scales, rotates that offset by the parent's rotation, and adds it to the mean. Child scales are divided by \(0.8N\) before mapping back through the inverse scale activation. Then the parent is removed. The samples are random in the parent's anisotropic local coordinates; the reference does not deterministically split along the single largest axis.
\end{itemize}

The cited source performs clone first and split second. It then prunes primitives whose opacity is below the scheduled minimum; after the configured size gate it also prunes those whose recorded screen radius is too large or whose maximum world scale exceeds \(0.1\) times the scene extent. It finally resets scheduled opacities by capping decoded opacity at \(0.01\), mapping back to logits, and resetting that parameter group's Adam moments. Those are choices of this revision, not universal constants or a claim that every fork follows them. (\href{https://github.com/graphdeco-inria/gaussian-splatting/blob/54c035f7834b564019656c3e3fcc3646292f727d/scene/gaussian_model.py\#L409-L472}{\texttt{densify\_and\_prune}, split and clone}; \href{https://github.com/graphdeco-inria/gaussian-splatting/blob/54c035f7834b564019656c3e3fcc3646292f727d/scene/gaussian_model.py\#L258-L261}{\texttt{reset\_opacity}})

The training-loop ordering matters. At the cited revision, \texttt{loss.backward()} runs first. Inside the no-gradient control block, the program accumulates visible gradient/radius statistics, performs any scheduled clone/split/prune, performs any scheduled opacity reset (also at the first density-control iteration for a white background), and only then calls the exposure and Gaussian optimizer steps and clears gradients. It does \textbf{not} do optimizer-step-then-density-edit as a generic pseudocode might suggest (\href{https://github.com/graphdeco-inria/gaussian-splatting/blob/54c035f7834b564019656c3e3fcc3646292f727d/train.py\#L142-L186}{\texttt{train.py}, lines 142 and 163--186}).

The model code must keep variable-length tensors and optimizer state aligned. When it appends children, it appends zero rows to Adam's \texttt{exp\_avg} and \texttt{exp\_avg\_sq}; when pruning, it slices those moments by the same keep-mask; opacity reset replaces the opacity parameter and zeros that parameter's two moment tensors (\href{https://github.com/graphdeco-inria/gaussian-splatting/blob/54c035f7834b564019656c3e3fcc3646292f727d/scene/gaussian_model.py\#L316-L405}{\texttt{GaussianModel}, lines 316--405}). This revision constructs new \texttt{Parameter} objects during append/prune/reset and does not transfer the just-backpropagated \texttt{.grad} tensor to those replacements. Consequently, on an edit step those replaced groups have no gradient for the immediately following optimizer call; their stored Adam moments are preserved/sliced or reset as the code specifies, and the next render/backward populates gradients for the new parameters. This behavior follows from the pinned PyTorch code path; it is an implementation fact to verify again when changing frameworks or revisions.

A small branch trace (numbers are teaching values, not source defaults): average projected-mean gradients are \([0.02,0.31,0.24]\), largest world scales are \([0.08,0.70,0.04]\), and the example gradient threshold is \(0.20\) with a small/large cutoff of \(0.10\). The first splat stays; the second splits; the third clones. If the second has opacity \(0.003\) and the example opacity floor is \(0.005\), its children may be created by the clone/split operation and then removed by the later prune mask if they inherit that low opacity. A faithful trace must apply the implementation's actual sequence, rather than silently moving the opacity prune ahead of growth. For each append, initialize child moments to zero; for parent removal, slice all parameter and per-splat statistic arrays with the same keep-mask. This example teaches the branch and ordering, not expected quality or useful production thresholds.

Represent the optimizer state and the topology schedule in a checkpoint, not just the visible means and colors. A reproducible experiment records calibrated cameras, image preparation, loss weights, optimizer type and learning-rate schedule, random seed, density thresholds and intervals, opacity-reset schedule, source revision, held-out views, memory use, and render/training time. The pinned loop tests held-out cameras when configured, but the selection and metric must be recorded; a test view used repeatedly to tune the model is no longer an untouched final evaluation.

\subsection{Reconstruction exercise: trace one fit and one topology edit}\label{04-visual-computing-ii-ch039-reconstruction-exercise-trace-one-fit-and-one-topology-edit}

\hypertarget{04-visual-computing-ii-ch039-reconstruction-exercise-trace-one-fit-and-one-topology-edit}{}

\begin{enumerate}
\def\labelenumi{\arabic{enumi}.}
\tightlist
\item
  Starting from the unit-footprint two-splat example, derive \(C\), \(\partial\ell/\partial c_1\), \(\partial\ell/\partial\alpha_1\), and \(\partial\ell/\partial a_1\). Then repeat the opacity-logit derivative with \(G_1=0.4\) while keeping \(o_1=0.5\) and the other quantities fixed. State why \(\alpha_1(1-\alpha_1)\) is not the general chain factor.
\item
  Trace the pinned clone/split decision for three splats with the values above. Add a low-opacity flag on the large splat; explain whether the reference grows first or prunes first and what happens to its children.
\item
  For each changed tensor, state what happens to its Adam moments and to the just-computed gradient. Name the exact code boundaries you would inspect after changing the pinned revision.
\end{enumerate}

\textbf{Worked answer and solution criteria.} The first trace is \(C=0.425\), \(\partial\ell/\partial c_1=0.0125\), \(\partial\ell/\partial\alpha_1=0.01875\), and at \(G_1=1,o_1=0.5\), \(\partial\ell/\partial a_1=0.0046875\). At \(G_1=0.4,o_1=0.5\), the effective opacity is \(\alpha_1=0.2\) and the logit chain factor is \(G_1o_1(1-o_1)=0.1\), not \(\alpha_1(1-\alpha_1)=0.16\). Recomputing the whole pixel with \(G_1=0.4\) gives \(\alpha_1=0.2\), \(C=0.2\), residual \(-0.2\), and \(\partial\ell/\partial\alpha_1=-0.15\); therefore \(\partial\ell/\partial a_1=-0.015\). For topology, use the accumulated-gradient threshold and scale cutoff to classify, perform clone/split in the pinned order, then apply opacity/size pruning. Appended Adam moment rows are zero; pruned rows are removed from moments and per-splat statistics; opacity reset zeroes that parameter's moments. Replaced parameter objects do not receive the previous \texttt{.grad} in this code. A complete answer distinguishes those source facts from any alternative policy, and it reports held-out image fit separately from geometry, memory, and runtime.

\subsection{A complete scalar training-step laboratory}\label{04-visual-computing-ii-ch039-a-complete-scalar-training-step-laboratory}

\hypertarget{04-visual-computing-ii-ch039-a-complete-scalar-training-step-laboratory}{}

The following complete Python program checks the two-splat color/logit derivatives against 240 finite differences, rejects the tempting wrong footprint chain rule, verifies a small SGD step, and tests stable identity and optimizer-moment alignment during a declared prune/clone policy. Run it with \texttt{python3\ tools/labs/gaussian\_training\_step\_lab.py} or with \texttt{python3\ -O\ tools/labs/gaussian\_training\_step\_lab.py}. It models one scalar pixel with fixed footprint and ordering; it does not train geometry, reproduce Graphdeco topology ordering, or validate GPU gradients. Implementation explanations are beside the operations in the code.

\begin{verbatim}
#!/usr/bin/env python3
"""A scalar two-splat gradient and topology-state lab, not a scene trainer."""
from dataclasses import dataclass, replace
import math
import random


def require(condition, message):
    # Checks remain active under python -O; these are executable lesson claims.
    if not condition:
        raise AssertionError(message)


def sigmoid(x):
    # Branch to avoid overflowing exp for a large negative or positive logit.
    if x >= 0:
        return 1 / (1 + math.exp(-x))
    e = math.exp(x)
    return e / (1 + e)


def loss_and_gradient(parameters, footprints, target, background=0.0):
    # The trainable vector is (front color, back color, front logit, back logit).
    # Geometry, sorting order and the Gaussian footprint values are held fixed.
    c1, c2, a1, a2 = parameters
    g1, g2 = footprints
    if not all(math.isfinite(v) for v in (*parameters, *footprints, target, background)):
        raise ValueError('All scalar inputs must be finite')
    if not (0 <= g1 <= 1 and 0 <= g2 <= 1):
        raise ValueError('This unclamped teaching model requires footprints in [0,1]')
    o1, o2 = sigmoid(a1), sigmoid(a2)
    alpha1, alpha2 = o1*g1, o2*g2
    behind = alpha2*c2 + (1-alpha2)*background
    color = alpha1*c1 + (1-alpha1)*behind
    residual = color-target
    # Chain rule: pixel alpha = opacity * footprint. Its logit derivative is
    # g*o*(1-o), NOT alpha*(1-alpha), unless the footprint equals one.
    dcolor = (alpha1, (1-alpha1)*alpha2,
              (c1-behind)*g1*o1*(1-o1),
              (1-alpha1)*(c2-background)*g2*o2*(1-o2))
    return .5*residual*residual, tuple(residual*d for d in dcolor), color


def finite_difference(parameters, footprints, target, axis, step=1e-6):
    # Re-evaluate the forward loss twice. This oracle does not use the analytic
    # derivative, so a missing transmittance or footprint factor is detectable.
    if not math.isfinite(step) or step <= 0:
        raise ValueError('Finite-difference step must be finite and positive')
    lo, hi = list(parameters), list(parameters)
    lo[axis] -= step
    hi[axis] += step
    return (loss_and_gradient(hi, footprints, target)[0]
            - loss_and_gradient(lo, footprints, target)[0])/(2*step)


@dataclass(frozen=True)
class ScalarSplat:
    # Only one scalar parameter and its two Adam moments are represented here.
    # Real trainers maintain aligned state for every trainable parameter group.
    identity: int
    parameter: float
    first_moment: float
    second_moment: float


def prune_and_clone(records, prune_ids, clone_ids):
    # This lab chooses prune-before-clone, with pruning winning conflicts.
    # It is a declared teaching policy, not Graphdeco's exact control order.
    identities = [r.identity for r in records]
    if len(identities) != len(set(identities)):
        raise ValueError('Stable identities must be unique')
    if not (set(prune_ids) | set(clone_ids)) <= set(identities):
        raise ValueError('An edit refers to an absent identity')
    output = [r for r in records if r.identity not in prune_ids]
    next_id = max(identities, default=-1) + 1
    for parent in records:
        if parent.identity in clone_ids and parent.identity not in prune_ids:
            # The child copies the parameter, but starts with zero moments.
            # Survivors retain their own moments despite changing array indices.
            output.append(replace(parent, identity=next_id,
                                  first_moment=0.0, second_moment=0.0))
            next_id += 1
    return output


def run():
    parameters = (.8, .2, 0., math.log(1/3))
    loss, gradient, color = loss_and_gradient(parameters, (1., 1.), .4)
    require(math.isclose(color, .425), 'Worked pixel color')
    require(math.isclose(gradient[0], .0125), 'Front color gradient')
    require(math.isclose(gradient[2], .0046875), 'Front logit gradient')
    # A small simultaneous update lowers this particular local loss. It does
    # not establish global convergence or identify a 3-D scene from one pixel.
    updated = tuple(p-.1*g for p, g in zip(parameters, gradient))
    require(loss_and_gradient(updated, (1., 1.), .4)[0] < loss, 'Local SGD descent')

    rng = random.Random(731)
    for _ in range(60):
        p = (rng.uniform(-.2, 1.2), rng.uniform(-.2, 1.2),
             rng.uniform(-3, 3), rng.uniform(-3, 3))
        footprints = (rng.uniform(.05, 1), rng.uniform(.05, 1))
        target = rng.uniform(-.1, 1.1)
        _, analytic, _ = loss_and_gradient(p, footprints, target)
        for axis in range(4):
            numerical = finite_difference(p, footprints, target, axis)
            require(math.isclose(analytic[axis], numerical, rel_tol=2e-5, abs_tol=1e-8),
                    'Finite-difference gradient disagreement')

    # A deliberate wrong formula must disagree when footprint != 1.
    _, correct, c = loss_and_gradient(parameters, (.35, .8), .4)
    alpha = sigmoid(parameters[2])*.35
    behind = sigmoid(parameters[3])*.8*parameters[1]
    wrong = (c-.4)*(parameters[0]-behind)*alpha*(1-alpha)
    require(not math.isclose(correct[2], wrong, rel_tol=1e-3),
            'The negative chain-rule example must expose the missing footprint')
    require(sigmoid(-1000.) == 0. and sigmoid(1000.) == 1., 'Finite saturation')

    original = [ScalarSplat(7, .2, .3, .4), ScalarSplat(9, .8, .7, .6)]
    edited = prune_and_clone(original, {7}, {7, 9})
    require([r.identity for r in edited] == [9, 10], 'Prune wins; child gets new identity')
    require(edited[0] == original[1], 'Survivor state follows identity, not old index')
    require(edited[1].parameter == .8 and edited[1].first_moment == 0
            and edited[1].second_moment == 0, 'Child state initialization')
    require(len(original) == 2, 'Topology edit leaves source checkpoint immutable')
    for bad in ([original[0], original[0]],):
        try:
            prune_and_clone(bad, set(), set())
        except ValueError:
            pass
        else:
            raise AssertionError('Duplicate identity accepted')
    try:
        finite_difference(parameters, (1, 1), .4, 0, step=0)
    except ValueError:
        pass
    else:
        raise AssertionError('Zero finite-difference step accepted')
    print('PASS: worked gradient, SGD, 240 finite differences, wrong-chain rejection, '
          'saturation, topology/moment identity and negative inputs; scalar CPU model only')


if __name__ == '__main__':
    run()
\end{verbatim}

\begin{center}\rule{0.5\linewidth}{0.5pt}\end{center}

\hyperref[volume-reading-map]{Reading map} · \hyperref[04-visual-computing-ii-ch038-projecting-ellipsoids-into-screen-space-conics]{Previous chapter} · \href{../04-Light-Transport-and-Image-Synthesis-P01/main.pdf\#nameddest=04-visual-computing-ii-ch040-lighting-reconstruction-and-surface-motion}{Next chapter (Light Transport and Image Synthesis, Part 1, Ch. 10)}
